\documentclass[conference]{IEEEtran}
% IEEEtran conference mode defaults to US Letter. Add a4paper only if the
% final submission venue explicitly requires A4.
\IEEEoverridecommandlockouts

\usepackage{cite}
\usepackage{amsmath,amssymb,amsfonts}
\usepackage{graphicx}
\usepackage{booktabs}
\usepackage{array}
\usepackage{textcomp}
\usepackage{xcolor}
\usepackage{url}
\usepackage{placeins}

% On-board result fields.
\newcommand{\oursLatency}    {51.0}
\newcommand{\oursThroughput} {20}
\newcommand{\oursPower}      {48.5}
\newcommand{\oursEfficiency} {0.43}
\newcommand{\oursPerplexity} {9.53}

\graphicspath{{./}}
\setlength{\textfloatsep}{7pt plus 2pt minus 2pt}
\setlength{\floatsep}{6pt plus 2pt minus 2pt}
\setlength{\intextsep}{7pt plus 2pt minus 2pt}
\setlength{\abovecaptionskip}{3pt}
\setlength{\belowcaptionskip}{0pt}

\begin{document}

\title{Accelerating LLaMA-2 Decoding on FPGAs via DDR-Bank-Aware Attention and Signed DSP Packing}

% Replace this controlled anonymous block with the final author information
% if the target venue is not double blind.
\author{\IEEEauthorblockN{Anonymous Authors}}

\maketitle

\begin{abstract}
Batch-one large-language-model decoding is strongly constrained by external-memory traffic, as projection weights must be repeatedly streamed while the key–value (KV) cache grows with sequence length. To address these bottlenecks, this paper presents a four-SLR FPGA accelerator for LLaMA2-7B decoding on the AMD–Xilinx Alveo U250. First, a signed W4A15 linear engine packs two low-bit weight products sharing the same activation into one DSP48E2, improving arithmetic utilization for memory-bound GEMV operations. Second, the architecture assigns one processing element and one DDR bank to each SLR, keeping weight and KV traffic local while overlapping data transfer and computation. Third, a compressed KV8 attention engine performs a single-pass online-softmax scan without materializing the complete attention-score vector. The routed design operates at 300 MHz and achieves a measured token latency of 51.0 ms at sequence position 0 with 48.5 W active board power. The W4A15-KV8 configuration achieves 9.53 perplexity on WikiText-2, matching the quantized software reference.

\end{abstract}

\begin{IEEEkeywords}
large language model, FPGA, INT4, DSP packing, KV cache, online softmax,
multi-SLR dataflow
\end{IEEEkeywords}

\section{Introduction}

Transformer models underpin both encoder-oriented language representations and
decoder-only generative LLMs~\cite{transformer,bert,llama2}. During batch-one
generation, decoder-only models execute all Transformer layers for every new
token. Their projections are matrix--vector multiplications (GEMVs) with low
arithmetic intensity, so external-memory traffic commonly limits performance
rather than peak arithmetic throughput~\cite{llama2,inference-roofline}. The
attention cost also increases with context length because each new query scans
the previously stored keys and values. Therefore, batch-one decoding is
strongly constrained by external-memory traffic.

FPGAs can specialize numerical formats, data paths, buffers, and memory layout
for this workload. On a multi-SLR FPGA, wide cross-SLR broadcasts also make
timing closure difficult. The design must keep high-bandwidth data local to each
SLR whenever possible.

We therefore choose the numerical format and memory mapping to match the U250
DSP48E2 and DDR organization. The contributions are:

\begin{itemize}
  \setlength{\itemsep}{1pt}
  \item \textbf{Signed W4A15 DSP packing.}
  Two W4 products sharing one A15 activation are accumulated in one DSP48E2
  and recovered exactly after G32 accumulation.

  \item \textbf{SLR-local execution with dedicated DDR-bank assignment.}
  Each SLR contains one processing element assigned to one DDR bank. Activation
  groups and RoPE coefficients are distributed once and then reused locally.

  \item \textbf{KV8 online attention.}
  Compressed K and V data are processed in a single pass using online softmax
  without storing the complete attention-score vector.
\end{itemize}

The accelerator targets batch-one LLaMA2-7B decoding on an Alveo U250. The reported
latency includes FPGA execution from a 4096-element input hidden-state vector
through all 32 decoder layers, the final RMSNorm, and the 4096-to-32000
vocabulary projection. It excludes tokenization, embedding lookup, sampling,
and host-side sequence control.

Each pre-normalized decoder layer contains RoPE-based
self-attention~\cite{rope}, a SwiGLU
feed-forward network, and two residual updates. Input token IDs and embedding
lookup are shown for model context, but hardware evaluation begins at the
embedded hidden state.

\begin{figure}[!htbp]
  \centering
  \includegraphics[width=0.78\columnwidth]{prism-uploads/Transformer.png}
  \caption{LLaMA2-7B decoder workload. The FPGA evaluates the embedded hidden
  state through the repeated decoder, final RMSNorm, and vocabulary head.}
  \label{fig:decoder}
\end{figure}

\section{Background and Related Work}

\subsection{Low-Bit Quantization}

Prior quantization methods reduce LLM numerical precision through outlier-aware,
reconstruction-based, or activation-aware
approaches~\cite{llmint8,gptq,smoothquant,awq,atom,aptq}. APTQ additionally considers
attention-output error~\cite{aptq}. We apply APTQ offline while
constraining weights to one symmetric W4 scale per $128\times256$ hardware tile.
Runtime projection inputs use signed A15 quantization in 32-element groups for
direct mapping to DSP48E2 arithmetic.

\subsection{Attention and FPGA Acceleration}

FlashAttention reduces GPU traffic through IO-aware tiling and fusion, but its
large query-block reuse is limited in batch-one decoding~\cite{flashattention}.
A stable online normalizer instead permits a single KV-cache scan without
storing the complete score vector. SwiftKV applies online normalization to an
edge-oriented attention accelerator~\cite{swiftkv}.

FlightLLM provides a mapping flow that accounts for compressed weights and
activations in FPGA LLM inference and uses multiple compute cores placed within
individual SLRs on the Alveo U280~\cite{flightllm}. QLlama combines MXINT8
quantization with dedicated
vector-matrix multiplication, attention, quantization, and basic-operator units
on the U280~\cite{qllama}. FlightOPU couples compute cores to external-memory
channels and uses SLR-aware placement and synchronization on a multi-die V80
FPGA~\cite{flightopu}. The U250 used here has four DDR banks and four SLRs. Our
design uses W4 weights with one scale per $128\times256$ tile, online A15
activations, signed DSP48E2 packing, and compressed KV-cache data stored in each
processing element's assigned DDR bank.

\section{Hardware Architecture}

\begin{figure*}[!t]
  \centering
  \includegraphics[width=0.85\textwidth]{prism-uploads/PACKED_W4A15_G32_EVALUATION.png}
  \caption{Conventional and packed W4A15 datapaths. The packed path accumulates
  a signed 46-bit result over one 32-element group before extracting the two
  packed dot products and applying a one-bit correction when the lower sum is
  negative.}
  \label{fig:packed}
\end{figure*}

\subsection{Overall System Organization}

The design is specialized for LLaMA2-7B: 32 layers,
residual width 4096, 32 heads of width 128, feed-forward width 11008,
vocabulary size 32000, and maximum context length 4096. Eight execution modes
select the Q, K, V, O, gate, up, down, and vocabulary projections.

\subsubsection{Processing Elements Placed by SLR with DDR-Bank Assignment}

Fig.~\ref{fig:four-slr} shows the physical organization of the accelerator.
Processing element $\mathrm{PE}_i$ is constrained to $\mathrm{SLR}_i$ and
communicates with DDR bank $i$ through its assigned 512-bit AXI master. Each
bank stores the round-robin output tiles assigned to that processing element,
the corresponding weight scales, KV-cache data for eight attention
heads, and striped output or residual data.

Within each processing element, weight transfers and KV-cache transfers use
the same AXI master connected to the assigned DDR bank. A single global
scheduler sequences all four processing elements and schedules these transfers.

\begin{figure}[!htbp]
  \centering
  \includegraphics[width=\columnwidth]{prism-uploads/arc_top.jpg}
  \caption{Four-SLR organization. Thick paths represent data streams and
  dashed paths represent the compact commands sent by the global scheduler to
  the four processing elements. Activation groups and RoPE coefficients are
  reused from local buffers after distribution. Weight and KV-cache accesses
  use each processing element's assigned DDR bank.}
  \label{fig:four-slr}
\end{figure}

The global scheduler sends a registered command to each processing element
through a dedicated FIFO. Linear commands contain 14 bits. Attention commands
contain 18 bits: a 6-bit layer index and a 12-bit token position. Each
processing element uses local command registers for address generation, so the
complete controller state is not broadcast across the four SLRs. The floorplan
does not assign the scheduler FSM to a specific SLR; it assigns the shared
activation and RoPE buffers to SLR1. For cross-SLR distribution, one 32-element
activation group is packed into a 512-bit packet containing 32 signed A15
values, or 480 payload bits, and one 32-bit FP32 activation scale. Each
processing element stores this activation packet in local replay memory and
reuses it across its local row blocks and column tiles.

At initialization, the complete 4096-position RoPE coefficient table is
preloaded into URAM. For each token
position, the corresponding RoPE coefficients are distributed to local caches
in each processing element. Because the same positional coefficients are
used by all decoder layers at a given token position, each processing element
reuses the cached coefficients for its eight assigned heads throughout the 32
layers.

\subsubsection{Four-Bank Memory Organization}

The architecture exposes four independent 512-bit AXI memory interfaces
(\texttt{gmem0}--\texttt{gmem3}) corresponding to the four DDR banks of the
target U250 platform. Weight tiles are assigned to the banks in a round-robin
manner, while four processing elements operate concurrently on streams from
their assigned banks. At 300 MHz, the four interfaces provide an ideal
aggregate payload ceiling of 76.8 GB/s (71.53 GiB/s).

Partitioned BRAM and URAM retain intermediate tensors on chip. Cross-SLR
communication consists of narrow control commands, activation groups, RoPE
coefficients for the current token position, and outputs gathered from the four
processing elements.

\subsection{W4A15 Linear Engine}

\subsubsection{Signed Arithmetic Packing}

The packed datapath uses one DSP48E2 instead of two multipliers for each W4 pair
(Fig.~\ref{fig:packed}). With two row pairs evaluated concurrently across 32
activation lanes, the packed inner-product core uses 64 DSP48E2s instead of
128.

\textbf{Step 1 (Fabric Packing):}
Two signed INT4 weights ($w_h$ and $w_l$) that multiply the same signed A15
activation $x_k$ are first sign-extended and packed at a 23-bit spacing using
a registered fabric adder:
\begin{equation}
 w_{p,k} =
 \left(\mathrm{sext}_{27}(w_{h,k}) \ll 23\right)
 +\mathrm{sext}_{27}(w_{l,k}).
 \label{eq:wpack}
\end{equation}
Equation~\eqref{eq:wpack} is an arithmetic addition, not a bitwise
concatenation, and it inserts no unused zero bits between the two packed fields.
Consequently, a negative $w_{l,k}$ sign-extends and borrows from the upper
field. The asymmetric two's-complement code $-8$ is intentionally excluded so
that the packing bound uses the symmetric bound $|w|\le 7$.

\textbf{Step 2 (Packed Multiplication):}
The 27-bit packed operand $w_{p,k}$ and the 15-bit activation $x_k$,
sign-extended to 18 bits, are multiplied in one DSP48E2 operating in
$27\times18$ mode. The DSP48E2 produces a 45-bit product. Under the W4A15
bounds used here, the product magnitude is below $2^{40}$ and the subsequent
packed accumulation over one 32-element group fits in 46 signed bits.

\textbf{Step 3 (G32 Accumulation):}
The 45-bit packed products are accumulated directly across the 32 elements of
one group:
\begin{equation}
 A_p=\sum_{k=0}^{31} w_{p,k} x_k
     =2^{23}D_h+D_l .
 \label{eq:packedacc}
\end{equation}
The A15 bound is selected with the 32-element group size and symmetric W4 so
that the 32-element dot product for each packed weight field fits within 23
signed bits. The same 23-bit separation allows the two W4 operands to fit
within the 27-bit DSP48E2 multiplier input because
$|7\cdot2^{23}+7|<2^{26}$. The maximum lane magnitude is bounded by
$32\times7\times16383=3{,}669{,}792<2^{22}$.

\textbf{Step 4 (Extracting the Packed Dot Products and Correcting Sign Borrow):}
The two 32-element integer dot products are recovered from the 46-bit
accumulator using
\begin{align}
 D_l &= \mathrm{signed}(A_p[22:0]),\\
 D_{h,\mathrm{raw}} &= \mathrm{signed}(A_p[45:23]),\\
 D_h &= D_{h,\mathrm{raw}} + \left(D_l[22]\, ?\, 1 : 0\right).
 \label{eq:correction}
\end{align}
The one-bit addition of $D_l[22]$ corrects the borrow from the lower packed
field whenever $D_l$ is negative.

After extracting the two packed dot products, scaling is performed as
\begin{equation}
 y_r \leftarrow y_r +
 \mathrm{FP32}(D_{r,g,T})\,S_{x,g}\,
 \mathrm{FP32}(S_{w,T}).
 \label{eq:dequant}
\end{equation}
In Eq.~\eqref{eq:dequant}, the scaled integer result and the accumulation over
all eight 32-element groups in a 256-column tile, as well as over all column
tiles, use FP32. Therefore, the 23-bit, 32-element integer dot products are not
accumulated across the complete matrix width. FP32 outputs are retained for
residual paths or requantized by Eq.~\eqref{eq:aquant} before a following
linear projection.

\subsubsection{Weight and Activation Quantization}

Each projection matrix is partitioned into $128\times256$ hardware tiles. For a tile
$T$, maximum-magnitude scaling provides the initial symmetric W4 scale:
\begin{equation}
\begin{aligned}
 S^{(0)}_{w,T}&=\frac{\max_{i\in T}|W_i|}{7},\\
 q_T(w;S)&=\mathrm{clip}_{[-7,7]}
 \left(\mathrm{round}(w/S)\right).
\end{aligned}
 \label{eq:tilequant}
\end{equation}
APTQ then refines the quantized weights offline while retaining one symmetric W4
scale per $128\times256$ tile and values restricted to $[-7,7]$. The W4 weights
are calibrated offline and stored with their FP16 scales.

Before each linear projection, the input activation vector is quantized online
into independent groups of 32 elements:
\begin{equation}
\begin{aligned}
 S_{x,g}&=\frac{\max_{k\in g}|x_k|}{16383},\\
 \hat x_k&=\mathrm{clip}_{[-16383,16383]}
 \left(\mathrm{round}\left(x_k/S_{x,g}\right)\right).
\end{aligned}
 \label{eq:aquant}
\end{equation}
The group scale $S_{x,g}$ remains FP32. These signed A15 activations are then
combined with streamed W4 weights inside the packed W4A15 datapath.

Each 512-bit word contains 32 INT4 weights from each of four output rows,
exactly filling the interface width. The 256 weight words of each tile are
stored contiguously, while FP16 scales are stored separately.

\subsubsection{Linear Dataflow and Load--Compute Overlap}

As illustrated in Fig.~\ref{fig:pipeline-overlap}, each linear processing
element overlaps weight transfer, packed integer computation, and output
generation through FIFO-decoupled pipeline stages. While the DDR reader streams
the weights of tile $(T+1)$, the packed W4A15 MAC processes tile $T$, and the
output stage completes tile $(T-1)$. After pipeline fill, these stages operate
concurrently on different tiles.

\begin{figure}[!ht]
  \centering
  \includegraphics[width=\columnwidth]{prism-uploads/overlap.png}
  \caption{Linear-processing-element load--compute overlap. FIFO-decoupled
  stages concurrently stream weights for tile $(T+1)$, execute the packed
  W4A15 MAC for tile $T$, and generate the dequantized output for tile $(T-1)$
  after pipeline fill.}
  \label{fig:pipeline-overlap}
\end{figure}
\FloatBarrier

\subsection{KV8 Online-Attention Engine}

\begin{figure}[!htbp]
  \centering
  \includegraphics[width=\columnwidth]{Attention.jpg}
  \caption{KV8 attention with one DDR bank assigned to each processing element.
  Each processing element scans eight attention heads from its assigned bank.
  The QK datapath, online-softmax recurrence, and parallel value-update datapath
  with four local accumulator banks process the compressed KV stream without
  storing the complete attention-score vector.}
  \label{fig:kv8}
\end{figure}

During autoregressive decoding, the KV8 attention engine scans the compressed
K/V records for each attention head once from the processing element's assigned
DDR bank. During the scan, it maintains the running maximum, normalization
term, and weighted value accumulator. The full score vector is never written to
memory. After the scan, the weighted accumulator is normalized to produce the
attention output.

\subsubsection{Compressed KV Block and DDR-Bank Head Assignment}

For one token and one attention head, K and V are quantized before being stored
in DDR. The compressed data occupy five 512-bit AXI transfers: one metadata
transfer, two transfers for K, and two transfers for V:
\[
 [\mathrm{metadata},K_0,K_1,V_0,V_1].
\]

Let $r$ be the signed integer encoding of a Q15.17 value in one K or V group,
and let $m$ be the maximum absolute encoded value in that group. The shift
amount and quantized value
are
\begin{align}
 q_s &= \max(\lfloor\log_2 m\rfloor-6,0),\\
 q &= \mathrm{sign}(r)\,
 \min\left(127,
 \left\lfloor\frac{|r|+2^{q_s-1}}{2^{q_s}}\right\rfloor\right).
 \label{eq:kvquant}
\end{align}
When $q_s=0$, the rounding offset is zero. K and V quantization is performed
online after RoPE because the stored values depend on the current token
position. The newly generated KV block is written once and forwarded
to the attention pipeline, avoiding an immediate DDR readback.

The K and V data for one token and one head occupy 320 bytes. FP16 K and V data
would occupy 512 bytes, so compression reduces the cache footprint by 37.5\%.

\subsubsection{Single-Pass Online Softmax}

The exponential is implemented with a lookup-based approximation, and final
normalization uses a reciprocal unit. The result is quantized to A15 before the
O projection.

\subsubsection{Parallel Value Update with Four Accumulator Banks}

Four parallel value-update units collectively process 16 elements per cycle.
The QK datapath has the same processing rate, so both require eight cycles for
one 128-element vector. All value-update state remains local to the processing
element.

RMSNorm, attention, and SwiGLU outputs are quantized to signed A15 values before
their following projections. Residual addition is integrated into the
writeback stages of the O and down projections.

\section{Experimental Evaluation}

\subsection{Experimental Setup}

The accelerator is implemented on the AMD--Xilinx Alveo U250 using Vitis HLS
and Vitis 2023.2 and operates at 300 MHz with four 512-bit AXI interfaces.
Board measurements use batch size one. Kernel latency is measured using device
profiling timestamps, excluding host-side tokenization, embedding lookup, and
sampling. Board power is measured using XRT sensors, and latency is averaged
over 100 runs.

\subsection{Post-Route Resource Results}

The packed inner-product core uses 64 DSP48E2s, compared with 128 for two
independent multipliers. A complete linear processing element uses 93 DSP48E2s.

Table~\ref{tab:resource} reports the post-route utilization of the implemented
U250 design. BRAM values are BRAM36 equivalents.

\begin{table}[!htbp]
\caption{Post-Route Resource Utilization on Alveo U250}
\label{tab:resource}
\centering
\scriptsize
\setlength{\tabcolsep}{8pt}
\begin{tabular}{lrr}
\toprule
Resource & Used & Utilization \\
\midrule
LUT       & 321.5k & 18.6\% \\
Flip-flop & 384.6k & 11.1\% \\
BRAM36    & 689    & 25.6\% \\
URAM      & 422    & 33.0\% \\
DSP48E2   & 846    & 6.9\% \\
\bottomrule
\end{tabular}
\end{table}

The routed design operates at 300 MHz.

\subsection{On-Board Performance and Power}

Table~\ref{tab:performance} reports the measured performance at sequence
position 0.

\begin{table}[!htbp]
\caption{Measured On-Board Performance at Sequence Position 0}
\label{tab:performance}
\centering
\scriptsize
\setlength{\tabcolsep}{3pt}
\begin{tabular}{lr}
\toprule
Metric & Value \\
\midrule
Kernel frequency & 300 MHz \\
Token latency & \oursLatency{} ms \\
Throughput & \oursThroughput{} token/s \\
Active board power & \oursPower{} W \\
Energy efficiency & \oursEfficiency{} token/J \\
\bottomrule
\end{tabular}
\end{table}

At sequence position 0, the measured latency is \oursLatency{} ms,
corresponding to approximately \oursThroughput{} token/s at \oursPower{} W. The reported
result is not sequence averaged.

\subsection{Model Quality and Agreement with the Software Reference}

Quality is evaluated with the same LLaMA2-7B checkpoint, WikiText-2 split,
tokenizer, context length, and sliding-window policy for the FP16 software
reference, the quantized software reference, and the FPGA output. Calibration
uses 128 sequences of length 2048 sampled from the C4 dataset independently of
the WikiText-2 evaluation set.

\begin{table}[!htbp]
\caption{Quality Evaluation on WikiText-2}
\label{tab:quality}
\centering
\scriptsize
\setlength{\tabcolsep}{3pt}
\begin{tabular}{lrrrr}
\toprule
Configuration      & W    & A    & KV   & PPL \\
\midrule
FP16 software      & FP16 & FP16 & FP16 & 9.39 \\
Quantized software & W4   & A15  & KV8  & 9.53 \\
FPGA               & W4   & A15  & KV8  & \oursPerplexity \\
\bottomrule
\end{tabular}
\end{table}

The FPGA achieves a WikiText-2 perplexity of \oursPerplexity{}, matching the
quantized software result.

\section{Conclusion}

We implemented a four-SLR LLaMA2-7B accelerator using signed W4A15 DSP packing
and compressed KV attention with dedicated DDR-bank assignment. The U250 design
operates at 300 MHz and achieves 20.83 token/s at sequence position 0 with
48.5 W board power. The W4, A15, and KV8 configuration obtains 9.53 WikiText-2
perplexity.

\begin{thebibliography}{00}

\bibitem{transformer}
A. Vaswani et al., ``Attention is all you need,'' in \emph{Proc. Adv. Neural
Inf. Process. Syst.}, vol.~30, pp. 5998--6008, 2017.

\bibitem{bert}
J. Devlin, M.-W. Chang, K. Lee, and K. Toutanova, ``BERT: Pre-training of deep
bidirectional transformers for language understanding,'' in \emph{Proc.
NAACL-HLT}, pp. 4171--4186, 2019.

\bibitem{llama2}
H. Touvron et al., ``Llama 2: Open foundation and fine-tuned chat models,''
arXiv:2307.09288, 2023.

\bibitem{inference-roofline}
Z. Yuan et al., ``LLM inference unveiled: Survey and roofline model
insights,'' arXiv:2402.16363, 2024.

\bibitem{llmint8}
T. Dettmers, M. Lewis, Y. Belkada, and L. Zettlemoyer, ``LLM.int8(): 8-bit matrix
multiplication for transformers at scale,'' in \emph{Proc. Adv. Neural Inf.
Process. Syst.}, vol.~35, pp. 30318--30332, 2022.

\bibitem{gptq}
E. Frantar, S. Ashkboos, T. Hoefler, and D. Alistarh, ``GPTQ: Accurate
post-training quantization for generative pre-trained transformers,'' in
\emph{Proc. Int. Conf. Learn. Representations (ICLR)}, 2023.

\bibitem{smoothquant}
G. Xiao, J. Lin, M. Seznec, H. Wu, J. Demouth, and S. Han, ``SmoothQuant:
Accurate and efficient post-training quantization for large language models,''
in \emph{Proc. Int. Conf. Mach. Learn. (ICML)}, pp. 38087--38099, 2023.

\bibitem{awq}
J. Lin et al., ``AWQ: Activation-aware weight quantization for on-device LLM
compression and acceleration,'' in \emph{Proc. Mach. Learn. Syst.}, vol.~6,
pp. 87--100, 2024.

\bibitem{atom}
Y. Zhao et al., ``Atom: Low-bit quantization for efficient and accurate LLM
serving,'' in \emph{Proc. Mach. Learn. Syst.}, vol.~6, pp. 196--209, 2024.

\bibitem{aptq}
Z. Guan, H. Huang, Y. Su, H. Huang, N. Wong, and H. Yu,
``APTQ: Attention-aware post-training mixed-precision quantization for large
language models,'' in \emph{Proc. 61st ACM/IEEE Design Automation Conf.
(DAC)}, 2024, doi: 10.1145/3649329.3658498.

\bibitem{flightllm}
S. Zeng et al., ``FlightLLM: Efficient large language model inference with a
complete mapping flow on FPGAs,'' in \emph{Proc. 2024 ACM/SIGDA Int. Symp.
Field-Programmable Gate Arrays (FPGA)}, 2024,
doi: 10.1145/3626202.3637562.

\bibitem{qllama}
H. Wen et al., ``QLlama: An FPGA-based microscaling quantization accelerator
for energy-efficient Llama2 inference,'' \emph{IEEE Embedded Syst. Lett.},
vol.~17, no.~5, pp. 337--340, 2025,
doi: 10.1109/LES.2025.3600563.

\bibitem{swiftkv}
J. Zhang et al., ``SwiftKV: An edge-oriented attention algorithm and
multi-head accelerator for fast, efficient LLM decoding,'' arXiv:2601.10953,
2026.

\bibitem{flightopu}
C. Wu et al., ``FlightOPU: An FPGA overlay processor for LLM with HBM-aware
multi-die architecture,'' in \emph{Proc. 2025 Int. Conf. Field Programmable
Technology (ICFPT)}, pp. 248--251, 2025,
doi: 10.1109/ICFPT67023.2025.00048.

\bibitem{flashattention}
T. Dao, D. Y. Fu, S. Ermon, A. Rudra, and C. R\'e,
``FlashAttention: Fast and memory-efficient exact attention with
IO-awareness,'' in \emph{Proc. Adv. Neural Inf. Process. Syst.}, vol.~35,
pp. 16344--16359, 2022.

\bibitem{rope}
J. Su, Y. Lu, S. Pan, A. Murtadha, B. Wen, and Y. Liu, ``RoFormer: Enhanced
Transformer with rotary position embedding,'' arXiv:2104.09864, 2021.

\end{thebibliography}

\end{document}
